% GENERATED FILE. Edit canonical lessons, then run npm run generate:books.
% canonical-source-hash: fnv1a64:d06c0f6bbbc1061c
% canonical-lessons: TA-C28-pirpakal, TA-C28-afternoon-boundary

\chapter{Afternoon}
\label{ch:ta-afternoon}

\begin{chapteropening}
\textbf{By the end of this chapter:} \emph{I can say \ta{பிற்பகல்} for afternoon, read it as \textquotedblleft{}after\textquotedblright{} plus \textquotedblleft{}daytime,\textquotedblright{} and name exactly where the dictionary's own noon and afternoon labels overlap.}

Show naṭuppakal straddling noon and afternoon, place matiyam as a Sanskrit noon-word, and label the unsourced moon claim unsupported.
\end{chapteropening}

\section[\ta{பிற்பகல்}]{\ta{பிற்பகல்} (piṟpakal) — \textquotedblleft{}afternoon,\textquotedblright{} a similar sandhi trick to \ta{நண்பகல்}}

\label{lesson:TA-C28-pirpakal}



\begin{quote}
You already met a sandhi-driven surface change in Chapter 17 — \ta{நள்} resurfacing as \ta{நண்}. This lesson's afternoon word does a similar kind of trick, with a different word — though be honest that the two changes aren't identical.
\end{quote}

\subsection*{You'll want to know: \ta{பிற்பகல்} — \ta{பின்} (after) + \ta{பகல்} (daytime), a similar sandhi story}
\textbf{\ta{பிற்பகல்}} (\textbf{piṟpakal}) — \textquotedblleft{}\textbf{afternoon}\textquotedblright{} — is a transparent compound: \textbf{\ta{பின்}} (\emph{piṉ}, \textquotedblleft{}\textbf{after}\textquotedblright{}) plus \textbf{\ta{பகல்}} (\emph{pakal}, \textquotedblleft{}\textbf{daytime},\textquotedblright{} already met inside \ta{நண்பகல்}, TA-C17 — though be aware the \ta{பிற்பகல்} Wiktionary page itself glosses \ta{பகல்} as \textquotedblleft{}noon\textquotedblright{} in its own etymology line, a small inconsistency across Wiktionary's own entries, not something this lesson introduces). Here's the familiar pattern: \ta{பின்} doesn't stay \ta{பின்} in this compound — it surfaces as \textbf{\ta{பிற்}-} before the following \ta{ப} — a \textbf{similar} kind of consonant sandhi to the one already taught in Chapter 17, where \textbf{\ta{நள்}} became \textbf{\ta{நண்}} before its own following \ta{ப}. Be honest, though: Wiktionary's own etymology for \ta{பிற்பகல்} doesn't discuss this mechanism at all, and the two changes move in different phonetic directions (nasal$\to$stop here, lateral$\to$nasal there) — this is the lesson's own observation, not a sourced equivalence. Wiktionary labels this word \textbf{formal register}.

\subsection*{Your turn}
Say these aloud:

\begin{itemize}
\raggedright
  \item \textquotedblleft{}piṟpakal\textquotedblright{} — \textquotedblleft{}afternoon,\textquotedblright{} piṉ (\textquotedblleft{}after\textquotedblright{}) + pakal (\textquotedblleft{}daytime\textquotedblright{})
  \item the sandhi echo — piṉ surfaces as piṟ-, a similar kind of change to TA-C17's naḷ$\to$naṇ, though not phonetically identical
\end{itemize}

\begin{tcolorbox}[breakable,colback=teal!4,colframe=teal!35!black,title={Before you move on}]
What does \ta{பிற்பகல்} literally build from, and what sandhi change happens to the first piece? (\textbf{\ta{பின்}} \textquotedblleft{}after\textquotedblright{} + \ta{பகல்} \textquotedblleft{}daytime\textquotedblright{} — \ta{பின்} surfaces as \ta{பிற்}-, a similar kind of sandhi to TA-C17's naḷ$\to$naṇ, though the two changes move in different phonetic directions, not identical mechanisms.) What register does the fetched page label? (\textbf{Formal}.) Next: inspect the real noon/afternoon boundary blur and the separate \emph{matiyam} root.
\end{tcolorbox}
\section[\ta{நடுப்பகல்} / \ta{மதியம்}]{Noon or afternoon? — let the sources keep their blur}

\label{lesson:TA-C28-afternoon-boundary}



\begin{quote}
\emph{Piṟpakal} transparently says after-daytime. Dictionaries still do not draw a perfectly hard semantic boundary around its synonyms.
\end{quote}

\subsection*{You'll want to know: \ta{நடுப்பகல்} overlaps}
Wiktionary's \emph{piṟpakal} page lists \textbf{\ta{நடுப்பகல்}} as an afternoon synonym, while the earlier noon lesson documents the same word as a synonym of \textbf{\ta{நண்பகல்}}, \textquotedblleft{}noon.\textquotedblright{} This is genuine noon/afternoon boundary blur, not an inconsistency to erase. Time-period categories can widen differently across speakers and sources.

\begin{cousinweb}
\textbf{\ta{மதியம்}} (\emph{matiyam}) is a Sanskrit tatsama from \emph{madhya}, \textquotedblleft{}middle,\textquotedblright{} the same root used by Kannada and Telugu. The fetched entry gives it \textbf{noon}, a synonym of \emph{naṇpakal}, not a dedicated afternoon meaning.

An initial search summary claimed \emph{matiyam} once meant \textquotedblleft{}moon,\textquotedblright{} but that claim is absent from the actual fetched page. It is therefore \textbf{not asserted}. Source inspection outranks a search snippet.
\end{cousinweb}

\subsection*{Your turn}
\begin{itemize}
\raggedright
  \item \emph{Say it:} \emph{naṭuppakal} can overlap noon and afternoon
  \item \emph{Say it:} \emph{matiyam} — Sanskrit \emph{madhya}, \textquotedblleft{}middle,\textquotedblright{} listed as noon
  \item \emph{Label:} the moon claim unsupported
\end{itemize}

\begin{tcolorbox}[breakable,colback=teal!4,colframe=teal!35!black,title={Before you move on}]
Does the source draw a hard noon/afternoon line? (\textbf{No}.) What is \emph{matiyam}'s sourced root and sense? (Sanskrit \emph{\textbf{madhya}}, \textquotedblleft{}middle,\textquotedblright{} and \textbf{noon}.) Is \textquotedblleft{}originally moon\textquotedblright{} confirmed? (\textbf{No}.)
\end{tcolorbox}
